% ===========================================================================
% methods.tex -- Methodology for
%   "Nonlinear Manifold Reduced Models with Precomputed Weak Operators"
%
% EVERY equation in this file was checked line by line against the implemented
% code in this repository.  Where the previous submission
% (paper/old-submission/main.tex) and the code disagree, THE CODE IS FOLLOWED
% and the disagreement is recorded in a comment beginning "DISCREPANCY".
%
% Sources checked (paths relative to the worktree root):
%   experiments/separable-decoder/sep_common.py          (features, head, SeparableDecoder)
%   experiments/mr-burgers2d/engines.py                  (spatial, residual, modes, weak,
%                                                         build_rom, nnls_capped,
%                                                         build_gauss_cold, make_lm)
%   experiments/mr-burgers2d/accuracy_paths.py           (make_stationary_lm, make_rom,
%                                                         make_diagnostics)
%   experiments/multiresolution-poisson/core.py          (assemble, rom_query, dst_solve)
%   experiments/multiresolution-poisson/correction_core.py (prepare_correction, correction_query)
%   experiments/multiresolution-poisson/kernel_solver.py (make_lm_kernel, guarded_gj)
%   experiments/cost-to-tolerance/ctol_tol.py            (lm_tau_poisson)
%   experiments/mr-heat2d/heat_core.py                   (mode_matrix, cn_factor, make_lm,
%                                                         make_rollout, make_query)
%   experiments/mr-heat2d/cp_algebra_paths.py            (make_gram_lm, build)
%   experiments/mr-heat2d/run_pilot.py                   (assemble, exactness assertion)
%   worktrees/2026-08-29-b2d-tensor/experiments/separable-decoder/b2d_tensor_common.py
%                                                        (backward_diff_bank_2d, build_T,
%                                                         symmetrize, q_of, dq_dh)
%
% NO NUMBER IS TYPED BY HAND.  Numeric slots use \gen{...} and are filled by the
% generator scripts listed in figures/README.md.
% ===========================================================================

\section{Method}
\label{sec:method}

We reduce a discretised PDE by restricting its state to the image of a frozen
learned decoder and solving, in the latent coordinates alone, a least-squares
problem built from a small set of smooth weak tests. The construction has four
parts: the trial manifold (\S\ref{sec:method:manifold}), the weak residual and
the solver that is actually run (\S\ref{sec:method:weak}), the offline
assembly of the reduced operators (\S\ref{sec:method:precomp}), and the three
PDE instances (\S\ref{sec:method:poisson}--\S\ref{sec:method:burgers}).
\S\ref{sec:method:eq} describes empirical quadrature, which is used only for
the operator that is not exactly precomputable. \S\ref{sec:method:tuning}
states precisely which quantities may be varied at deployment, and
\S\ref{sec:method:limits} states what the construction does not provide.

\subsection{Discrete setting and notation}
\label{sec:method:setting}

All problems are posed on $[0,1]^2$ with homogeneous Dirichlet data and
discretised by second-order finite differences on a uniform grid of $N$
intervals per axis. The unknown is the vector of interior nodal values
$u\in\Real^{n}$, $n=(N-1)^2$; boundary values are structurally zero and are not
unknowns. We write $\Aop\in\Real^{n\times n}$ for the negative discrete
Laplacian (five-point stencil, ghost zeros on all four walls) and $f\in\Real^n$
for a source or right-hand side. Reduced quantities carry three sizes that must
be kept distinct: $R$ is the width of the spatial bank, $k$ the latent
dimension solved for online, and $M$ the number of retained weak tests. Where
sampling is used, $m$ is the number of quadrature nodes. The implementations
assert $M\ge 4k$ and $m\ge 4M$ at build time, so the latent problem is always
overdetermined and the quadrature problem is never under-constrained.

\subsection{Trial manifold}
\label{sec:method:manifold}

The trial manifold is the image of an affine lift composed with a frozen
nonlinear coefficient map,
\begin{equation}
  \umanifold(\latent) \;=\; \lift \;+\; \bank\,\head(\latent),
  \qquad
  \bank\in\Real^{n\times R},\quad
  \head:\Real^{k}\to\Real^{R},\quad k\le R\ll n .
  \label{eq:decoder}
\end{equation}
Both $\lift$ and $\bank$ are fixed once the mesh is chosen, and $\theta$ is
frozen for every result reported here. In the implemented cells the Dirichlet
data are homogeneous, so $\lift=0$, and the boundary condition is enforced
structurally by a smooth vanishing factor folded into the bank itself:
\begin{equation}
  \bank_{x,:} \;=\; \bcmask(x)\,g_\phi(x)\T,
  \qquad
  \bcmask(x) \;=\; 16\,x_1(1-x_1)\,x_2(1-x_2),
  \label{eq:bank}
\end{equation}
where $g_\phi:\Real^2\to\Real^{R}$ is a random-Fourier-feature
MLP~\citep{Tancik2020FourierFeatures} evaluated pointwise at the nodal
coordinates. Because $\bcmask$ vanishes on $\partial[0,1]^2$, every state in the
manifold satisfies the boundary condition exactly and
$\partial \umanifold/\partial\latent$ vanishes there as well, for every
$\latent$ and at every resolution.
% DISCREPANCY (D1) vs old-submission Eq. (1): the previous manuscript wrote
% \tilde u = m \odot D(z) + u_g with a BINARY mask m and a nonzero lift u_g.
% The implemented decoder (sep_common.features / bc_poly) uses the SMOOTH
% polynomial factor above, folded into the bank, with u_g = 0.  Enforcement is
% still exact, but by a different mechanism, and as implemented it covers only
% homogeneous Dirichlet data.  We follow the code.

The coefficient map is a two-layer SiLU MLP plus a linear skip,
\begin{equation}
  \head(\latent) \;=\; \varphi_\theta(\latent) \;+\; W\T\latent,
  \qquad W\in\Real^{k\times R},
  \label{eq:head}
\end{equation}
so that the decoder Jacobian
\begin{equation}
  \jac{\umanifold}(\latent) \;=\; \frac{\partial \umanifold}{\partial \latent}
  \;=\; \bank\,\dhead(\latent),
  \qquad
  \dhead(\latent) \;=\; D\varphi_\theta(\latent) + W\T ,
  \label{eq:jacobian}
\end{equation}
retains a latent-independent component. The rank statement this permits is
exactly the following, and no more:
\begin{equation}
  \rank \jac{\umanifold}(\latent)
  \;\le\; \min\{\,k,\;\rank \bank\,\}\;\le\;\min\{k,R\},
  \label{eq:rank-bound}
\end{equation}
with equality to $k$ if and only if $\dhead(\latent)$ has rank $k$ and
$\operatorname{range}\dhead(\latent)\cap\ker\bank=\{0\}$. The linear skip makes
the first condition generic but does not prove it: $D\varphi_\theta(\latent)$
can cancel $W\T$ at particular $\latent$. We therefore measure the rank of the
projected Jacobian at every solved query rather than asserting it
(\S\ref{sec:method:weak}), and report the retained bank rank from the singular
values of $\bank$.
% Code: the Poisson kernel (correction_core.correction_query) computes the
% singular values of the projected Jacobian and sets projected_jacobian_rank_valid
% per query; run_pilot.assemble records retained_bank_rank for the heat bank.

Two consequences of \eqref{eq:decoder} are used throughout. First, the manifold
lies inside the $R$-dimensional linear span of the bank, so the representation
error of the manifold is bounded below by the linear projection error of that
span: nonlinearity in $\head$ buys a smaller \emph{solved} dimension $k$, not
an escape from the bank's span. Second, all $x$-dependence factors through
$\bank$, so restricting the decoder to any point set $S$ is exactly the cached
matrix $\bank_{S,:}\in\Real^{|S|\times R}$ times $\head(\latent)$. This is the
property that makes both node sampling (\S\ref{sec:method:eq}) and exact
operator preassembly (\S\ref{sec:method:precomp}) available from the same
decoder.
% DISCREPANCY (D6) vs old-submission Sec. 4.4: the previous manuscript's
% architecture was a ViT encoder feeding a LinearCPDecoder (linear skip +
% shallow MLP + rank-R canonical-polyadic spatial head).  The implemented
% decoder here has NO ENCODER in the deployed path: latent codes come from an
% auto-decoder fit during training and from a least-squares fit at query time
% (Sec. 5.3).  The per-node evaluation property that the CP head was chosen for
% is inherited by the separable bank.  The "ViT for global coupling" argument
% of the old manuscript does not apply to the method described here.

\subsection{Weak residual, objective, and solver}
\label{sec:method:weak}

\paragraph{Tests and row scaling.}
Let $\tests\in\Real^{M\times n}$ collect $M$ smooth test vectors as rows,
together with the discrete integration weight. In every cell the tests are the
tensor-product sine vectors
$\sin(\pi a x_1)\sin(\pi b x_2)$ sampled at interior nodes, retained for the
$M$ smallest eigenvalues $\lambda_i$ of $\Aop$, and scaled by a fixed
normalisation: the Poisson and Burgers cells use the orthonormal discrete sine
basis ($\tests\tests\T=I_M$), while the heat cell uses continuum-normalised
tests multiplied by the cell area $N^{-2}$, i.e.\ $N^{-1}$ times the
orthonormal basis, a uniform scaling chosen so that entries are comparable
across meshes. These vectors are exactly the eigenvectors of the five-point
operator, which gives the identity we use throughout,
\begin{equation}
  \tests\Aop \;=\; \rowscale\,\tests,
  \qquad
  \rowscale \;=\; \diag(\lambda_1,\dots,\lambda_M),
  \qquad
  \lambda_{(a,b)} \;=\; 4N^2\!\left[\sin^2\!\tfrac{\pi a}{2N}+\sin^2\!\tfrac{\pi b}{2N}\right].
  \label{eq:eigen-identity}
\end{equation}
% Code: this identity is asserted numerically in run_pilot.assemble, which
% checks || low_lam * (Psi^T G h) - Psi^T (-Lap_h)(G h) || / ||.|| < 1e-11 for
% both the operator and its Jacobian.  It is the reason the linear part of
% every weak residual below needs no quadrature.
The residual rows are additionally scaled by a diagonal $\rowscale_\star^{-1}$
whose choice is stated per PDE below; the scaling is part of the objective, not
a cosmetic normalisation, and it changes which solution the least-squares
problem selects.

\paragraph{The objective that is minimised.}
Given a fully discrete residual $\rfull(\cdot)$ for the PDE at hand, the
reduced problem solved online is the row-scaled, weak-tested least-squares
problem
\begin{equation}
  \latent^{\star}
  \;=\;
  \argmin_{\latent\in\Real^{k}}
  \tfrac12\,\norm{\rweak(\latent)}_2^2,
  \qquad
  \rweak(\latent) \;=\; \rowscale_\star^{-1}\,\tests\,\rfull\!\big(\umanifold(\latent)\big)
  \;\in\;\Real^{M},
  \qquad M>k .
  \label{eq:objective}
\end{equation}
This is a least-squares Petrov--Galerkin condition with an explicitly chosen
test space and row norm, in the sense of \citet{Carlberg2011LSPG} and
\citet{LeeCarlberg2020}. It is \emph{not} tangent-space Galerkin: the test
space is $\tests$, fixed and independent of $\latent$, not
$\jac{\umanifold}(\latent)$; and the system is overdetermined
($M>k$) rather than square.
% DISCREPANCY (D2) vs old-submission Eqs. (2)-(3): the previous manuscript
% stated tangent Galerkin, r_red(z) = J_D(z)^T (K u(z) - F) = 0, with the update
% dz = -(J_D^T K J_D)^{-1} r_red.  No implemented path forms J_D^T K J_D, and
% none imposes a square tangent-projected condition.  We follow the code.
% Note also that for a LINEAR trial space the two conditions still differ:
% residual least squares imposes V^T A^T (AVz-b)=0 while Galerkin imposes
% V^T(AVz-b)=0.  The rebuttal's claim that they coincide for a linear decoder
% holds only under additional assumptions on the test space and norm, which we
% do not make.

\paragraph{The update that is taken.}
Write $\jac{}(\latent)=\partial\rweak/\partial\latent\in\Real^{M\times k}$,
evaluated by forward-mode automatic differentiation through $\head$
\citep{Bradbury2018JAX}, $H=\jac{}\T\jac{}$ and $g=\jac{}\T\rweak$. Each
attempt solves the damped normal system
\begin{equation}
  \big(H + \lambda\,\diag(\diag H)\big)\,\delta\latent \;=\; -\,g,
  \label{eq:lm}
\end{equation}
and accepts $\latent\leftarrow\latent+\delta\latent$ only if $\delta\latent$ is
finite, $\norm{\delta\latent}\le\Delta$ for a fixed trust radius $\Delta$, and
the residual norm strictly decreases. On acceptance
$\lambda\leftarrow\max(\lambda/3,10^{-12})$; on rejection the step is discarded
and $\lambda\leftarrow\min(10\lambda,\lambda_{\max})$. This is damped
Levenberg--Marquardt with a monotone accept/reject test
\citep{Marquardt1963}; there is no line search.
% DISCREPANCY (D8): the old manuscript calls this "Gauss-Newton" throughout and
% prints an undamped update.  Implemented: LM with diagonal damping,
% lambda_0 = 1e-6 (Poisson kernel_solver.make_lm_kernel and Burgers
% accuracy_paths.make_stationary_lm) or 1e-4 (heat heat_core.make_lm), and an
% explicit trust radius taken from the spread of the training codes.
The curvature dropped in \eqref{eq:lm} is exactly
$\sum_{i}(\rweak)_i\,\nabla^2_{\latent}(\rweak)_i$. Since every implemented
$\rweak$ is affine or quadratic in $\head(\latent)$ and $\head$ is the only
nonlinearity in $\latent$, this dropped term is the residual-weighted curvature
of the head, $\sum_i (\rweak)_i \sum_a B_{ia}\nabla^2_\latent \headc{a}$ in the
linear-operator case. We identify it here rather than absorbing it into the
name of the method.

\paragraph{Stopping.}
Three families of exit are recorded and never conflated. The primary one is the
scale-free stationarity measure
\begin{equation}
  \station{\latent}{\cdot}
  \;=\;
  \frac{\norm{\jac{}(\latent)\T\rweak(\latent)}_2}
       {\normF{\jac{}(\latent)}\;\norm{\rweak(\latent)}_2},
  \qquad\text{exit when } \station{\latent}{\cdot}\le\eta_{\mathrm{tol}} ,
  \label{eq:stationarity}
\end{equation}
with $\normF{\cdot}$ the Frobenius norm. The second is a residual threshold:
relative, $\norm{\rweak(\latent)}\le\tau\,\norm{\rweak(\latent_0)}$, in the
elliptic cell, and absolute, scaled by the norm of the fitted input, in the
Burgers time step. The third family collects \emph{stalls}: exhausting the
attempt budget $I_{\max}$, a step below $10^{-14}(1+\norm{\latent})$, the
damping reaching $\lambda_{\max}$, or a non-finite value. Each solve returns its
exit reason, and outputs produced under a stall exit are reported as
early-stopped outputs, never as converged solves.
% Code: reason codes differ per cell and we report them per cell rather than
% unifying them.  Burgers (accuracy_paths.make_stationary_lm): 4 stationarity,
% 1 residual, 2 tiny step, 3 rejected at lambda >= 1e14, 0 budget.  Poisson
% (kernel_solver.make_lm_kernel): 6 stationarity, 2 relative residual tau,
% 1 stall, 3 damping limit, 5 non-finite initial value, 0 budget.  Heat
% (heat_core.make_lm): 1 stationarity, 2 tiny step, 3 damping limit,
% 4 non-finite, 0 budget.  The heat criterion divides by ||J||_F only, applied
% to a residual already normalised by ||target||; Poisson and Burgers use the
% ratio in Eq. (8) directly.

\paragraph{What stationarity does and does not certify.}
$\station{\latent}{\cdot}\le\eta_{\mathrm{tol}}$ certifies that $\latent$ is a
first-order critical point of the specific objective \eqref{eq:objective} on
the frozen manifold. It does not certify a global minimum of that objective; it
says nothing about the $n-M$ residual components that no test vector sees; it
does not bound the representation error
$\min_{\latent}\norm{\umanifold(\latent)-u^{\star}}$ of the frozen decoder; and
it therefore does not bound the physical error of the returned field. We
report weak stationarity and physical error as separate quantities throughout,
and we record cases in which the weak residual decreases while the field error
does not.

\paragraph{Initialisation.}
No query uses the solution it is predicting. The elliptic cell starts from the
cached training code whose stored projected weak prediction is nearest the
projected source. The heat cell takes the better of the nearest library code and
the mean code after a QR-compressed full-field least-squares fit. The Burgers
cell fits the latent code to the supplied initial field on a fixed
$48\times 48$ Gauss--Legendre product rule that does not depend on the query
mesh: with $Q\mathcal{R}$ the thin QR factorisation of the weight-scaled bank
evaluated at the Gauss nodes and $\tilde u$ the bilinear interpolation of the
dense input at those nodes, the fit minimises
$\norm{\mathcal{R}\,\head(\latent)-Q\T\tilde u}$ by the same LM iteration, started
at the training code minimising the corresponding quadratic surrogate. The
interpolation of the supplied input is charged to the query.

\subsection{Precomputed weak operators}
\label{sec:method:precomp}

\paragraph{Linear operators.}
For a fixed linear discrete operator $\Aop$, substituting \eqref{eq:decoder}
into $\rfull(u)=\Aop u-f$ and applying $\tests$ gives an exactly precomputable
reduced form,
\begin{equation}
  \tests\big(\Aop\,\umanifold(\latent)-f\big)
  \;=\; B\,\head(\latent) \;-\; b,
  \qquad
  B \;=\; \tests\Aop\bank\in\Real^{M\times R},
  \qquad
  b \;=\; \tests\,(f-\Aop\lift)\in\Real^{M}.
  \label{eq:linear-precomp}
\end{equation}
In the implemented cells the sine tests diagonalise $\Aop$, so
$B=\rowscale\,\tests\bank$ by \eqref{eq:eigen-identity} and, after row scaling
by $\rowscale^{-1}$, the reduced matrix is simply $\tests\bank$: a separable
sine transform of the bank, assembled offline by two one-dimensional
contractions.
% Code: core.assemble builds B with einsum('xa,xyr,yb->abr') from the skinny
% sine factor S; run_pilot.assemble builds it as Psi^T G.  No quadrature fit and
% no sampling is involved in either.

\paragraph{Polynomial operators.}
For a field operator that is polynomial in $u$ with fixed coefficients, the
projected term is a fixed multilinear form in $\head$. For the quadratic
advection used here, with $DG$ the bank differenced by the fixed stencil,
\begin{equation}
  \tests\,\big[\,\bank\head \odot (D^-_{x}+D^-_{y})\bank\head\,\big]_{i}
  \;=\;
  \sum_{a,b} T_{iab}\,\headc{a}\headc{b},
  \qquad
  T_{iab}=\sum_{x}\tests_{ix}\,\bank_{xa}\,(DG)_{xb},
  \label{eq:tensor}
\end{equation}
with $DG=D^-_{x}\bank+D^-_{y}\bank$ and ghost zeros on all four walls. $T$ is
not symmetric in $(a,b)$ because its two factors are the bank and the
differenced bank, so we store the symmetrised table $Q=T+T^{(a\leftrightarrow
b)}$ and evaluate
\begin{equation}
  q(\head)=\tfrac12\,\head\T Q\,\head,
  \qquad
  \frac{\partial q}{\partial \head} = Q\,\head .
  \label{eq:tensor-eval}
\end{equation}
Storage is $MR^2$ entries and one residual evaluation costs $O(MR^2)$; the
matrix $Q\head$ formed for the residual is reused for the Jacobian, so a
residual-and-Jacobian pair costs $O(MR^2+MRk)$. Bank width therefore enters the
online cost quadratically and independently of the latent dimension. The offline
build costs $O(nMR^2)$ and is blocked over $x$ so that the $n\times R^2$
intermediate is never materialised.

\paragraph{Exactly when this is available.}
Three conditions must hold together. (i) The linear operator is fixed, or
affinely separable in its parameters, so that $B$ can be assembled once and
reused; an arbitrary new coefficient field does not inherit the same
offline/online split. (ii) The field operator is polynomial in $u$ with fixed
coefficients. (iii) \emph{For the Burgers advection specifically, the
sign-dependent upwind switch must be inactive}: the stencil reduces to the fixed
backward difference on both axes only where $u>0$, so \eqref{eq:tensor} is an
identity only if the state is non-negative at every stencil node involved. This
must hold not only for the data and not only for the converged decoded state,
but for \emph{every trial state the solver visits}. Non-negative training data is
not sufficient; decoded states can undershoot below zero even when the truth
does not, and we report the measured fraction of negative decoded values and the
resulting departure from the sign-upwind operator
(\gen{positivity audit: fraction of negative decoded nodes and tensor-vs-upwind
parity per mesh}) beside every exactness statement. We do not change the
full-order discretisation to make the algebra polynomial.

\paragraph{What is and is not mesh-independent.}
Assembly of $\bank$, $\tests$, $B$ and $T$ is mesh-dependent and must be redone
at each resolution even when the network weights $\theta$ transfer unchanged;
we charge and report it separately as setup. The cached reduced evaluation --
$B\head$, $Q\head$, the $k\times k$ damped solve \eqref{eq:lm} -- contains no
explicit loop over mesh nodes, so its operation count depends on $(M,R,k)$
alone. This does not imply a mesh-independent complete query: reading a dense
input and writing a dense requested field must touch $O(n)$ values, and the
dense decode $\bank\head(\latent)$ costs $O(nR)$. Nor does it imply
mesh-independent iteration counts or mesh-independent accuracy. We measure
complete device-query cost with its input, setup, solve and output components
reported separately.

\subsection{Elliptic instance: Poisson}
\label{sec:method:poisson}

For $\Aop u=f$ we take $\rowscale_\star=\rowscale$, so by
\eqref{eq:eigen-identity} and \eqref{eq:linear-precomp} the weak residual is
\begin{equation}
  \rweak(\latent)
  \;=\;\rowscale^{-1}\tests(\Aop\umanifold(\latent)-f)
  \;=\; B_0\,\head(\latent) - b_0,
  \qquad
  B_0=\tests\bank,\quad b_0=\rowscale^{-1}\tests f .
  \label{eq:poisson-residual}
\end{equation}
Because the retained tests are orthonormal and $\Aop u^{\star}=f$, this residual
is exactly the tested error,
\begin{equation}
  \rweak(\latent)\;=\;\tests\big(\umanifold(\latent)-u^{\star}\big),
  \label{eq:poisson-error-identity}
\end{equation}
so the minimised objective \eqref{eq:objective} is the squared discrete $L^2$
error of the reduced state in the retained modes, not merely a residual norm.
This is the precise sense in which the elliptic solve minimises an error
(energy-type) functional; it is a consequence of the row scaling and holds only
for the modes $\tests$ retains.
% Code: core.assemble puts the scaling in `project` (W = lam^{-1}) and leaves B
% unscaled, which is algebraically identical to Eq. (13).
% DISCREPANCY (D2b) vs old-submission Eq. (3): the previous manuscript described
% this solve as Gauss-Newton on a tangent-projected system without identifying
% its objective.  The objective is Eq. (13)-(14).

\paragraph{Analytically eliminated linear corrections.}
The deployed elliptic model augments the frozen coefficient vector with a small
number $c$ of fixed linear directions $C\in\Real^{R\times c}$ and solves
\begin{equation}
  \min_{\latent\in\Real^k,\;y\in\Real^{c}}
  \; \norm{B_0\big(\head(\latent)+Cy\big)-b_0}_2^2 .
  \label{eq:poisson-augmented}
\end{equation}
Let $B_0C=Q\mathcal{R}$ be a thin QR factorisation. Since $Q$ does not depend on
$\latent$, the inner minimisation over $y$ has the closed-form solution
$\mathcal{R}y=Q\T\big(b_0-B_0\head(\latent)\big)$ and the minimised value is
\begin{equation}
  \min_{y}\norm{B_0(\head(\latent)+Cy)-b_0}
  \;=\;
  \norm{B_\perp \head(\latent)-b_\perp},
  \qquad
  B_\perp=(I-QQ\T)B_0,\quad b_\perp=(I-QQ\T)b_0 .
  \label{eq:varpro}
\end{equation}
So \eqref{eq:poisson-augmented} is solved by running the LM iteration of
\S\ref{sec:method:weak} on $\latent$ alone with the deflated operator, then
recovering $y$ with one triangular solve. This is variable projection
\citep{GolubPereyra1973} with a \emph{constant} projector, which matters for the
derivative bookkeeping: because $I-QQ\T$ is independent of $\latent$, the
elimination introduces no approximation at all, and
$\jac{}(\latent)=B_\perp\,\dhead(\latent)$ is the \emph{exact} Jacobian of the
eliminated residual. The only derivative term dropped anywhere in the elliptic
solve is the Gauss--Newton head curvature identified after \eqref{eq:lm}; the
$y$ block contributes none, because it enters linearly and is eliminated in
closed form. We report both the reduced stationarity of the deflated problem
and the augmented stationarity of \eqref{eq:poisson-augmented}, whose gradient
is $\big[\jac{}\T \rfull;\,(B_0C)\T \rfull\big]$ normalised by
$\big(\normF{\jac{}}^2+\normF{B_0C}^2\big)^{1/2}\norm{\rfull}$, together with the
backward error of the triangular recovery and the numerical rank of $B_0C$.

\subsection{Parabolic instance: heat}
\label{sec:method:heat}

The semi-discrete heat equation is $\dot u=-\kappa\Aop u$. The implementation
advances it with Crank--Nicolson,
\begin{equation}
  \Big(I+\tfrac{\Delta t\,\kappa}{2}\Aop\Big)u^{n+1}
  \;=\;
  \Big(I-\tfrac{\Delta t\,\kappa}{2}\Aop\Big)u^{n},
  \label{eq:cn}
\end{equation}
% DISCREPANCY (D3a): the old manuscript states BACKWARD EULER (its Eq. (4)).
% The implemented scheme is Crank-Nicolson (heat_core.cn_factor).  We follow the code.
and the reduced step is obtained by substituting the manifold into the
\emph{fully discrete} equation \emph{before} projecting. With
$u^{n}=\umanifold(\latent_n)$ the fully discrete residual is
\begin{equation}
  \rfull_n(\latent)
  \;=\;
  \Big(I+\tfrac{\Delta t\kappa}{2}\Aop\Big)\umanifold(\latent)
  \;-\;
  \Big(I-\tfrac{\Delta t\kappa}{2}\Aop\Big)\umanifold(\latent_n),
  \label{eq:heat-full-residual}
\end{equation}
which is nonlinear in $\latent$ because $\umanifold$ is. Projecting with
$\tests$, using \eqref{eq:eigen-identity}, and row-scaling by
$\rowscale_\star=I+\tfrac{\Delta t\kappa}{2}\rowscale$ gives the implemented
residual
\begin{equation}
  \rweakn{n}(\latent)
  \;=\;
  B_0\,\head(\latent) \;-\; D\,B_0\,\head(\latent_n),
  \qquad
  D=\diag\!\left(\frac{1-\tfrac{\Delta t\kappa}{2}\lambda_i}
                      {1+\tfrac{\Delta t\kappa}{2}\lambda_i}\right),
  \qquad B_0=\tests\bank,
  \label{eq:heat-step}
\end{equation}
and the step is defined as
$\latent_{n+1}=\argmin_{\latent}\norm{\rweakn{n}(\latent)}_2$, solved by the LM
iteration of \S\ref{sec:method:weak} warm-started at $\latent_n$. Two points
are structural and were both wrong in the previous presentation.
% DISCREPANCY (D3b) vs old-submission Eq. (4): the previous manuscript printed
%   (J^T J + dt kappa J^T K J) z_{n+1} = J^T J z_n,
% a LINEAR system in z_{n+1} obtained by replacing \tilde u(z) with J_D z.  That
% substitution is valid only for a homogeneous linear decoder: it discards the
% decoder bias, the current operating point, and the linearisation remainder.
% The implemented step is the nonlinear least-squares problem above.
% DISCREPANCY (D9): the old manuscript also wrote a latent ODE
% J^T J \dot z = -kappa J^T K \tilde u(z).  No implemented path forms a latent
% mass matrix J^T J or integrates an ODE in z; the step is a per-step
% minimisation.
First, the right-hand side of \eqref{eq:heat-step} is built from the
\emph{decoded current state} $\head(\latent_n)$ at every step. Carrying instead
the quantity that the previous step already made small freezes the recursion:
the rollout then reproduces a one-step solution to round-off and the error is
silently wrong. Second, $\latent_{n+1}$ appears through $\head$, so the step is
a nonlinear least-squares problem and not a linear solve.

\paragraph{Cholesky and contracted normal equations.}
Because $B_0$ and $\dhead$ are small, the damped normal matrix in \eqref{eq:lm}
is symmetric positive definite by construction, and the deployed path
(\texttt{nmrom\_cholesky}) factors it by Cholesky and applies two triangular
solves instead of a general LU factorisation. A further variant avoids
assembling the $M\times k$ weak Jacobian at all: with the fixed Gram matrix
$S=B_0\T B_0\in\Real^{R\times R}$ assembled offline,
\begin{equation}
  H=\dhead\T S\,\dhead,
  \qquad
  g = \dhead\T\Big(S\big(\head(\latent)-\head(\latent_0)\big)+B_0\T\rweakn{n}(\latent_0)\Big),
  \label{eq:gram}
\end{equation}
where anchoring the gradient at the initial iterate avoids cancellation near a
good warm start. This replaces per-iteration $O(MR)$ work by $O(R^2)$ and is
exact; we report it as an algebraic ablation of the same step, not a different
model.

\subsection{Nonlinear advection: Burgers}
\label{sec:method:burgers}

The full-order problem is $u_t+u(u_x+u_y)=\nu\Delta u$ with the non-conservative
sign-upwind stencil
\begin{equation}
  N(u)_{ij} \;=\; u_{ij}\,\big(\delta_x u + \delta_y u\big)_{ij},
  \qquad
  (\delta_x u)_{ij} = N\!\cdot\!
  \begin{cases}
    u_{ij}-u_{i-1,j}, & u_{ij}>0,\\[2pt]
    u_{i+1,j}-u_{ij}, & u_{ij}\le 0,
  \end{cases}
  \label{eq:upwind}
\end{equation}
with $\delta_y$ defined analogously on the second index and switching on the
\emph{same} centre value $u_{ij}$,
with ghost zeros on all four walls, and
backward Euler in time,
$\rfull_n(u)=u-u^{n}+\Delta t\big(N(u)-\nu\Delta_h u\big)$. Projecting and row
scaling by $\rowscale_\star=I+\Delta t\,\nu\rowscale$ gives
\begin{equation}
  \rweakn{n}(\latent)
  \;=\;
  \big(I+\Delta t\,\nu\rowscale\big)^{-1}
  \Big[\,
    B_0\head(\latent)-B_0\head(\latent_n)
    \;+\;\Delta t\,\big(\,\widehat{N}(\latent)+\nu\rowscale B_0\head(\latent)\,\big)
  \Big],
  \label{eq:burgers-step}
\end{equation}
where $B_0=\tests\bank$ and the diffusion term is exact by
\eqref{eq:eigen-identity}. The row scaling is precisely the diagonal of the
linearised implicit operator in the modal basis, i.e.\ the modal form of the
Helmholtz preconditioner used by the full-order solver, so the reduced and
full-order problems are scaled consistently.

Only the advection projection $\widehat{N}(\latent)\approx\tests N(\bank\head(\latent))$
requires a choice. Two are implemented and compared:
\begin{align}
  \text{sampled (EQ)}\quad
  &\widehat{N}(\latent)=\sum_{i\in\mathcal S} w_i\,\tests_{:,i}\;N_i\big(\bank\head(\latent)\big),
  && \text{sign switch retained,}
  \label{eq:burgers-eq}\\[2pt]
  \text{precomputed}\quad
  &\widehat{N}(\latent)=\tfrac12\,\head(\latent)\T Q\,\head(\latent),
  && \text{exact iff $u\ge 0$ at every stencil node.}
  \label{eq:burgers-tensor}
\end{align}
The sampled form \eqref{eq:burgers-eq} evaluates the decoder only at the
five-point stencil of each retained node, which is the cached
$5m\times R$ block times $\head(\latent)$, and it reproduces the sign-dependent
stencil exactly at those nodes. The precomputed form
\eqref{eq:burgers-tensor} removes sampling entirely but carries the sign
restriction of \S\ref{sec:method:precomp} and departs from
\eqref{eq:upwind} wherever the decoded or trial state undershoots below zero.
Time stepping uses a fixed substep count per output interval; each step is
warm-started at $\latent_n$, and a linear extrapolation
$\latent_n+(\latent_n-\latent_{n-1})$ is used instead when, and only when, its
residual norm is smaller.

\subsection{Empirical quadrature}
\label{sec:method:eq}

The sampled rule of \eqref{eq:burgers-eq} is fitted \emph{offline}, from decoder
outputs, by non-negative least squares \citep{LawsonHanson1974NNLS} in the
empirical-quadrature form of \citet{Hernandez2017EQ} and
\citet{YanoPatera2019LPEQ}. A candidate node set $\mathcal C$ is drawn once with
a fixed seed. For each of $n_{\text{fit}}$ stored training codes $\latent^{(\ell)}$ we
evaluate the decoded state, form the per-node advection contributions
$N_i(\bank\head(\latent^{(\ell)}))$ and the exact full-grid target
$\tests N(\bank\head(\latent^{(\ell)}))$, and stack
\begin{equation}
  \min_{w\ge 0}\;
  \norm{\,\mathcal{G}\,w-\beta\,}_2,
  \qquad
  \mathcal{G}_{(\ell,i),\,c}=\tests_{i,c}\,N_c\big(\bank\head(\latent^{(\ell)})\big),
  \qquad
  \beta_{(\ell,i)}=\big[\tests N(\bank\head(\latent^{(\ell)}))\big]_i ,
  \label{eq:nnls}
\end{equation}
with the rows normalised by their Euclidean norms before the fit. The support is
grown greedily by the Lawson--Hanson criterion with an exact non-negative refit
after each addition and a hard cap of $m$ nodes; if the support saturates below
$m$, it is padded with the remaining candidates of largest mean design magnitude
and refit, so the achieved node count equals the requested one. Three counts that
the previous manuscript did not separate are separated here: the \emph{requested}
node count $m$, the \emph{achieved} support size, and the number of \emph{fit
snapshots} $n_{\text{fit}}$ used to assemble \eqref{eq:nnls}.
% DISCREPANCY (D4) vs old-submission Sec. 4.3: the previous manuscript fitted EQ
% against the TANGENT-PROJECTED residual J_D^T r with k rows per snapshot, and
% presented the sample count as a deployment-time knob.  Implemented: the fit
% targets the WEAK ADVECTION TERM only, with M rows per snapshot; the linear
% terms are exact and never sampled; and selecting a different m selects a
% different STORED rule.  The old manuscript's irreconcilable counts (minimum
% support 300, n_eq in {4,...,64}, N_eq in [512,2560]) are the three quantities
% separated above.

A deployed rule is a stored triple: node indices, non-negative weights, and the
cached $m\times 5\times R$ bank block at the five-point stencil of each node.
Changing $m$ requires re-solving \eqref{eq:nnls} offline. At deployment one
selects among rules that already exist; nothing is re-optimised online. The
number of quadrature nodes is therefore a \emph{discrete selection among
precomputed rules}, not a continuous accuracy knob.

Two further cautions apply. The NNLS relative fit attained on its own snapshots
is recorded, but it is a property of the fitting set and is not a fidelity
certificate on unseen states; we certify a rule by measured full-grid parity and
by physical error, never by its NNLS residual. And the selected support is one
admissible solution among several: rebuilding the same problem under a different
floating-point summation order can select different nodes.

\paragraph{What per-node decoder evaluation saves, and what it does not.}
Because all $x$-dependence factors through $\bank$, evaluating the decoder at
the $5m$ stencil coordinates costs $O(mR)$ and is independent of $n$: this is
the property the previous CP spatial head was chosen for, and the separable bank
inherits it. It does not reduce the cost of the dense input projection, of the
dense output decode $\bank\head(\latent)$ ($O(nR)$), or of the requested
full-field output; and it buys nothing for the linear terms, which are already
exact and reduced by \eqref{eq:linear-precomp}.

\subsection{Deployment-time controls at a fixed checkpoint}
\label{sec:method:tuning}

All controls below act on a single frozen artefact. The network weights
$\theta$, the bank $\bank$ and its width $R$, the latent dimension $k$, the test
set $\tests$ and its size $M$, the time discretisation, and the query contract
are identical across every operating point we report. What may change is:
\begin{enumerate}\itemsep2pt
  \item the LM attempt budget $I_{\max}$, per solve for the elliptic problem and
        per time step for the evolution problems;
  \item the stopping tolerance --- the stationarity threshold
        $\eta_{\mathrm{tol}}$ of \eqref{eq:stationarity}, and where applicable the
        residual threshold, relative ($\tau$) for the elliptic cell and absolute
        for the Burgers step;
  \item the index of the stored empirical-quadrature rule, where sampled
        evaluation is used at all.
\end{enumerate}
Nothing else is varied, and nothing is re-optimised at query time. An operating
point whose solves exited on budget or on a stall is labelled as such and its
outputs are reported as early-stopped outputs; losses are reported as losses. We
\emph{measure} the resulting error-versus-cost relation. We do not assume it is
monotone or that it forms a frontier: a smaller weak residual need not produce a
smaller field error, and no amount of additional optimisation can pass the
representation limit of the frozen decoder, which we report separately as the
best achievable fit of the bank and of the head.
% DISCREPANCY (D5) vs old-submission Sec. 5.1: the previous manuscript fed the
% projected operator to a conjugate-gradient solver, raising a symmetry question
% once row weights are introduced.  No implemented path applies CG to the
% projected operator: the k x k damped normal system is solved directly, by
% Gauss-Jordan elimination with symmetric diagonal scaling and a backward-error
% gate with a charged general-solve fallback (kernel_solver.guarded_gj), or by
% Cholesky (heat).  The symmetry objection does not arise, but only because the
% implementation differs from the printed description.

\subsection{Intrusive and non-intrusive comparison}
\label{sec:method:intrusive}

The model described here is intrusive. It requires the discrete operator
$\Aop$, the stencil, and the right-hand side at inference, and its offline stage
requires assembling $B$ and, where used, $T$ at the deployment resolution. Neural
operators such as FNO \citep{Li2020FNO} and DeepONet \citep{Lu2021DeepONet} are
non-intrusive: they map data to data and apply where no discretisation or solver
is available at all. The two are therefore not substitutes, and a head-to-head
accuracy-per-cost table is meaningful only under one stated condition, which we
state before any comparison: a discretisation exists and may be used. Under that
condition the intrusive route can check and reduce its own equation error, while
the non-intrusive route cannot; outside it, the comparison does not apply. We
also hold the training data and the train/validation split fixed between the
compared models and record the index hashes, so that differences are not
attributable to different data.

\subsection{What the construction does not provide}
\label{sec:method:limits}

We state the boundaries of the method with it, rather than in a closing
paragraph.
\begin{itemize}\itemsep2pt
  \item \emph{No accuracy certificate.} Weak stationarity
        \eqref{eq:stationarity} is a first-order condition on $M$ tested
        components of a residual with $n\gg M$ components. It bounds neither the
        untested components nor the physical error.
  \item \emph{No escape from the bank.} By \eqref{eq:decoder} the manifold lies
        in $\operatorname{range}\bank$; the bank's linear projection error is a
        floor on the achievable field error, whatever the solver does.
  \item \emph{Offline cost is mesh-dependent.} Frozen weights transfer across
        meshes, but $\bank$, $\tests$, $B$ and $T$ do not; rebuilding them is
        setup work that we charge and report separately.
  \item \emph{Exact preassembly is conditional.} It requires a fixed or affinely
        separable linear operator and a polynomial field operator, and for the
        Burgers upwind stencil it additionally requires the sign restriction of
        \S\ref{sec:method:precomp}, which decoded and trial states can violate.
  \item \emph{Parameter-varying operators need a declared treatment.} A finite
        affine decomposition can be preassembled term by term; an arbitrary new
        coefficient field cannot inherit the same split.
  \item \emph{No convergence theory.} We have no convergence guarantee for the
        cold-start LM iteration, and no bound on the empirical-quadrature error
        as a function of $m$ on unseen instances. Both are stated as open.
  \item \emph{Classical solvers remain strong on the separable cells.} For
        constant-coefficient Poisson and heat on a rectangle, the discrete sine
        transform gives a direct solve \citep{Swarztrauber1977}; we name it as
        the faster comparator wherever it applies and report it alongside every
        iterative comparator.
\end{itemize}
